\section*{Chapter \ref{ch:b1:fluid-statics} --- Fluid Statics}

\begin{solution}{exo:b1:fluid-statics:1}
$P = P_0 + \rho gh$: \qty{2.0e5}{Pa} (\qty{2.0}{bar}), \qty{1.1e6}{Pa}
(\qty{11}{bar}), \qty{1.0e7}{Pa} (\qty{101}{bar}). At \qty{300}{m}: $P =
\qty{3.1e6}{Pa}$, force \qty{3.1}{MN} on \qty{1}{m^2} (net, against
\qty{1}{atm} inside: \qty{3.0}{MN}).
\end{solution}

\begin{solution}{exo:b1:fluid-statics:2}
$h = P/\rho g$: mercury \qty{0.760}{m}; water \qty{10.3}{m}. A pump can at
best create a vacuum above the water: the atmosphere then pushes the
column up by \qty{10.3}{m}, no more.
\end{solution}

\begin{solution}{exo:b1:fluid-statics:3}
$F_2 = F_1S_2/S_1 = \qty{10}{kN}$ (one tonne); $d_2 = d_1S_1/S_2 = \qty{2.0}{mm}$
(same volume displaced, same work).
\end{solution}

\begin{solution}{exo:b1:fluid-statics:4}
Same pressure at the oil--water interface level in both arms: $\rho_og
h_o = \rho_wgh_w$, $h_w = 0.85 \times 12 = \qty{10.2}{cm}$ of water above
that level in the other arm; the oil surface stands $12 - 10.2 = \qty{1.8}{cm}$
higher.
\end{solution}

\begin{solution}{exo:b1:fluid-statics:5}
$H = RT/Mg = 8.314 \times 273/(0.029 \times 9.81) = \qty{8.0}{km}$; $P/P_0 =
\eu^{-z/H}$: $0.54$ at \qty{5}{km}, $0.33$ at \qty{8848}{m}; half at $H\ln2 =
\qty{5.5}{km}$ --- the numbers used in \cref{ch:b1:kinetic-theory}.
\end{solution}

\begin{solution}{exo:b1:fluid-statics:6}
$F = \tfrac12\rho gh^2L = \qty{1.96e8}{N}$; \omterm{prop:b1:fluid-statics:wall}{center of pressure} at \qty{13.3}{m}
depth, i.e.\ \qty{6.7}{m} above the base; moment about the base $F \times
6.7 = \qty{1.3e9}{N.m}$. Curving the dam toward the water turns the thrust
into compression of the arch, carried to the valley sides.
\end{solution}

\begin{solution}{exo:b1:fluid-statics:7}
$917/1025 = 89\%$; log $60\%$. Raft: buoyancy at full immersion $1000 \times
2.0 \times 9.81 = \qty{19.6}{kN}$, own weight \qty{11.8}{kN}: \qty{7.8}{kN}
spare, ten people... the tenth with wet feet: nine to stay dry
(\qty{800}{kg}).
\end{solution}

\begin{solution}{exo:b1:fluid-statics:8}
Immersed volume $m/\rho$: water \qty{20}{cm^3}; brine $20/1.1 = \qty{18.2}{cm^3}$,
$\qty{1.8}{cm^3}$ less, i.e.\ $1.8/0.50 = \qty{3.6}{cm}$ more stem: \qty{7.6}{cm}
emerge. Sensitivity $\approx m/(\rho^2S) \approx \qty{40}{cm}$ per unit of relative
density (\qty{3.6}{cm} per $0.1$).
\end{solution}

\begin{solution}{exo:b1:fluid-statics:9}
Displaced volume $2800/1.025 = \qty{2732}{m^3}$: \qty{268}{m^3} ($9\%$)
emerge. Neutral: mass $1025 \times 3000 = \qty{3075}{t}$: take in \qty{275}{t}.
Fresh water: buoyancy drops to \qty{3000}{t}: \qty{75}{t} heavy --- it sinks
unless \qty{75}{t} are pumped out.
\end{solution}

\begin{solution}{exo:b1:fluid-statics:10}
Air displaced: $10^{-2} \times 1.20 = \qty{12.0}{g}$; helium inside $12.0 \times
4/29 = \qty{1.7}{g}$; rubber \qty{3.0}{g}: net lift \qty{7.3}{g} --- seven meters
of light string.
\end{solution}

\begin{solution}{exo:b1:fluid-statics:11}
Moment $\int_0^h\rho gx\cdot xL\,\dd x = \rho gLh^3/3$, force $\rho gLh^2/2$:
arm $2h/3$. From $h_1$ to $h_2$: moment $\rho gL(h_2^3 - h_1^3)/3$, force
$\rho gL(h_2^2 - h_1^2)/2$: depth $\frac23(h_2^3 - h_1^3)/(h_2^2 - h_1^2)$.
\end{solution}

\begin{solution}{exo:b1:fluid-statics:12}
Extra immersion $x$ adds the buoyancy $\rho gSx$ upward: $m\ddot x = -\rho gSx$,
$\omega_0^2 = \rho gS/m$, $T = 2\pi\sqrt{m/\rho gS} = 2\pi\sqrt{0.020/(1000 \times
9.81 \times 5\times10^{-5})} = \qty{1.3}{s}$.
\end{solution}

\begin{solution}{pb:b1:fluid-statics:1}
\textbf{1.} $P = 1.013\times10^5 + 1025 \times 9.81 \times 300 = \qty{3.12e6}{Pa}
= \qty{31}{bar}$.

\textbf{2.} $S = \pi(0.40)^2 = \qty{0.503}{m^2}$: net force $(P - P_0)S =
\qty{1.52e6}{N}$ --- 155 tonnes, a loaded truck on a hatch.

\textbf{3.} Keel \qty{10}{m} deeper: $+\rho g \times 10 = \qty{1.0e5}{Pa}$ more,
$3\%$: the hull is loaded almost uniformly.

\textbf{4.} $\Delta\rho/\rho = \chi_T\Delta P = 4.6\times10^{-10} \times 3.0\times10^6
= 0.14\%$: negligible for the pressure (a few kilopascals out of three
million).

\textbf{5.} Thirty bars outside, one inside: the hull is a pressure
vessel loaded inward; a cylinder (or sphere) turns the pressure into
compression along its wall, which steel resists well, whereas a flat
plate would bend.

\textbf{6.} $\Delta P\,R/e = 3.0\times10^6 \times 5.0/0.030 = \qty{5.0e8}{Pa}
= \qty{500}{MPa}$: close to the yield stress --- \qty{300}{m} is near the
limit for this hull, and deeper boats need thicker or stronger steel
(or titanium).

\textbf{7.} Displaced volume $m_0/\rho = 2800/1.025 = \qty{2732}{m^3}$: \qty{268}{m^3}
emerge ($9\%$).

\textbf{8.} Neutral: $m = \rho V = \qty{3075}{t}$: take in \qty{275}{t}.

\textbf{9.} Buoyancy falls by $0.1\%$ of $\rho Vg$: $0.001 \times 3075 \times 9.81
\times10^3 = \qty{30}{kN}$ downward (three tonnes). Deeper $\Rightarrow$ more
compression $\Rightarrow$ heavier: unstable --- a submarine must trim
continuously (and a steel hull compresses less than water, which helps;
a too-flexible hull would sink without return).

\textbf{10.} Buoyancy $1000 \times 3000 \times 9.81 = \qty{29.4}{MN}$ against
weight $3075 \times 9.81 \times 10^3 = \qty{30.2}{MN}$: \qty{0.74}{MN} down
(\qty{75}{t}); pump out \qty{75}{t}.

\textbf{11.} Immediately \qty{20}{t} light: take in \qty{20}{t} of water
(\qty{19.5}{m^3}) into compensating tanks.

\textbf{12.} Water can only be expelled by something at higher pressure
than the sea outside; compressed air in bottles pushes it out --- only
as long as the bottle pressure exceeds the ambient pressure at that
depth.

\textbf{13.} $PV$ constant: \qty{275}{m^3} at \qty{31}{bar} needs $275 \times 31
= \qty{8500}{m^3}$ of air at \qty{1}{atm} --- stored at \qty{200}{bar} in
\qty{43}{m^3} of bottles.

\textbf{14.} \qty{2.0}{bar} at \qty{10}{m}; \qty{4.0}{bar} at \qty{30}{m}.

\textbf{15.} $PV$ constant: $6.0/4.0 = \qty{1.5}{L}$ --- the ribcage is
crushed in; free divers feel it.

\textbf{16.} Air taken at \qty{4}{bar} expands fourfold: \qty{24}{L} in lungs
that hold \qty{6}{L}: rupture. Never hold your breath on ascent; exhale
continuously.

\textbf{17.} $\Delta P = \rho g \times 1.0 = \qty{1.0e4}{Pa}$; force $\qty{1.0e4}{Pa}
\times 0.05 = \qty{500}{N}$ on the chest: the breathing muscles cannot
lift fifty kilograms --- a snorkel works only in the first decimeters.

\textbf{18.} $12 \times 200 = \qty{2400}{L}$ at \qty{1}{bar}, i.e.\ \qty{600}{L}
at \qty{4}{bar}: \qty{30}{min}.

\textbf{19.} $V = 10^8/1025 = \qty{97600}{m^3}$.

\textbf{20.} In fresh water it must displace $10^5\ \unit{m^3}$: \qty{2400}{m^3}
more, i.e.\ $2400/8000 = \qty{0.30}{m}$ deeper.

\textbf{21.} $10^4/1.025 = \qty{9756}{m^3}$ more: \qty{1.2}{m} deeper.

\textbf{22.} Empty, the ship's $G$ sits high and little hull is immersed:
the metacenter may fall below $G$ and the ship capsize in a swell. Sea
water low in the hull lowers $G$ and deepens the draught.

\textbf{23.} $\rho gh = \qty{2.0e5}{Pa}$, \qty{0.2}{MN/m^2}: fifteen times less
than the submarine's \qty{3}{MN/m^2}.

\textbf{24.} No: the pressure at the submarine's depth is $P_0 + \rho gh$,
fixed by the depth alone. The tanker's weight is borne by the water it
displaces, which raises the sea level everywhere by an immeasurable
amount; no extra load reaches the submarine.

\textbf{25.} $\dd P/\dd z = -\rho g$ --- hence $P = P_0 + \rho gh$ for every
depth, force and lung volume here --- and Archimedes' theorem, which is
that relation integrated over a closed surface.
\end{solution}
